%%%%%%%%%%  scan idx 453  |  folio 440  %%%%%%%%%%
et la
suite exacte des \(\underline{\Tor}\) obtenue par tensorisation avec \(\cO_{X'}\):
\[
0\longrightarrow
\Tor_1^{\cO_X}(\cO_{X'},\cO_Y)
\longrightarrow
J\otimes_{\cO_X}\cO_{X'}
\longrightarrow
\cO_{X'}
\longrightarrow
\cO_{Y'}
\longrightarrow0.
\]
L'image du morphisme : \(J\otimes_{\cO_X}\cO_{X'}\to\cO_{X'}\) est
\(J.\cO_{X'}=\cO_{X'}(1)\). Par suite, on obtient la suite exacte:
\[
0\longrightarrow
\underline{\Tor}_1^{\cO_X}(\cO_{X'},\cO_Y)
\longrightarrow
J\otimes_{\cO_X}\cO_{X'}
\longrightarrow
\cO_{X'}(1)
\longrightarrow0.
\]

On fait alors son produit tensoriel par \(\cO_{Y'}\); comme
\(\underline{\Tor}_1^{\cO_X}(\cO_{X'},\cO_Y)\) est un \(\cO_{Y'}\)-Module, et que
\(\cO_{X'}(1)\) est inversible, on obtient la suite exacte:
\[
0\longrightarrow
\underline{\Tor}_1^{\cO_X}(\cO_{X'},\cO_Y)
\longrightarrow
g^*(N)
\longrightarrow
\cO_{Y'}(1)
\longrightarrow0.
\]
On vérifie aisément que le morphisme \(g^*(N)\to\cO_{Y'}(1)\) figurant dans cette suite est le morphisme canonique, de sorte qu'elle s'identifie à la suite (3.1.2). Par définition, on a donc:
\[
\underline{\Tor}_1^{\cO_X}(\cO_{X'},\cO_Y)=F.
\]

\textbf{3.3.} On suppose maintenant que \(X\) est un schéma ayant un faisceau ample. Alors il en est de même pour \(X'\) qui est projectif sur \(X\), et pour \(Y\) et \(Y'\). Par ailleurs, \(i\) et \(j\) sont des morphismes parfaits, puisque ce sont des immersions régulières; \(g\) est un morphisme lisse donc parfait car \(Y'=\bfP(N)\); enfin, \(f\) est parfait d'après 1.9. Ces morphismes étant projectifs, on peut définir les homomorphismes additifs (IV
\[
i_*:K^\circ(Y)\to K^\circ(X);\quad
j_*:K^\circ(Y')\to K^\circ(X');\quad
f_*:K^\circ(X')\to K^\circ(X);\quad
g_*:K^\circ(Y')\to K^\circ(Y).
\]

